% =============================================================================
% ANEXO Y — Eventos en Tiempo Real y Notificaciones
% =============================================================================
\chapter{Anexo Y — Eventos en Tiempo Real y Notificaciones}
\label{cap:anexo_realtime}

Catálogo de los canales de Supabase Realtime y de las notificaciones del sistema, base de las
funcionalidades en vivo del Sprint~3 (Capítulo~\ref{cap:sprint3}).\par

\section{Canales Realtime}
\label{sec:anexo_rt_canales}
\renewcommand{\arraystretch}{1.3}
\fontsize{9}{10.8}\selectfont\sffamily
\begin{tabularx}{\textwidth}{|l|X|l|}
\hline
\rowcolor{colorPrimario}
\textcolor{white}{\textbf{Canal}} & \textcolor{white}{\textbf{Evento difundido}} & \textcolor{white}{\textbf{Origen}} \\ \hline
Chat & Nuevo mensaje / borrado lógico. & \comando{chat\_messages} (V7) \\ \hline
\rowcolor{grisTecnico}
Presencia & Estado en línea/desconectado. & Canal de \textit{presence} \\ \hline
Likes & Nuevo \textit{like} al portafolio propio. & \comando{profile\_likes} (RLS V15) \\ \hline
\end{tabularx}\normalfont\normalsize
\par\vspace{0.3cm}

La difusión de \textit{likes} en vivo se habilitó con la política RLS de \comando{SELECT} de la
migración V15: sin ella, Supabase Realtime no emitía eventos al dueño del portafolio y el contador
no se actualizaba en tiempo real.

\section{Notificaciones del Sistema}
\label{sec:anexo_rt_notif}
\renewcommand{\arraystretch}{1.3}
\fontsize{9}{10.8}\selectfont\sffamily
\begin{tabularx}{\textwidth}{|l|X|}
\hline
\rowcolor{colorPrimario}
\textcolor{white}{\textbf{Endpoint}} & \textcolor{white}{\textbf{Función}} \\ \hline
GET \comando{/api/v1/notifications} & Lista las notificaciones del usuario. \\ \hline
\rowcolor{grisTecnico}
PATCH \comando{/api/v1/notifications/\{id\}/read} & Marca una notificación como leída. \\ \hline
POST \comando{/api/v1/notifications/read-all} & Marca todas como leídas. \\ \hline
\end{tabularx}\normalfont\normalsize

\begin{NotaTecnica}
En el frontend, las notificaciones se reflejan con \textit{toasts} (Sonner) y un contador en la
\comando{Navbar}; el \comando{notificationsStore} mantiene el estado sincronizado con los eventos
Realtime y con los endpoints REST.
\end{NotaTecnica}

\clearpage
